\documentclass[10pt]{extarticle}
\usepackage[a4paper,margin=1.5cm]{geometry}
\usepackage[T1]{fontenc}
\usepackage{amsmath,amssymb}
\usepackage{graphicx}
\usepackage[export]{adjustbox}
\usepackage{multicol}
\usepackage{enumitem}
\usepackage{booktabs,array,tabularx}
\setlength{\columnsep}{0.7cm}
\setlength{\parindent}{0pt}
\setlength{\parskip}{3pt}
\pagestyle{empty}
\makeatletter
\newcommand{\qfigdef}[4]{\@namedef{qf@file@#1}{#2}\@namedef{qf@w@#1}{#3}\@namedef{qf@wide@#1}{#4}}
\newcommand{\qf@check}[1]{\@ifundefined{qf@file@#1}{\PackageError{score2}{Unknown figure id #1}{}}{}}
\newcommand{\qf@img}[4][]{\includegraphics[width=#2\dimexpr\@nameuse{qf@w@#4}pt\relax,max width=#3,max height=0.45\textheight,#1]{\@nameuse{qf@file@#4}}}
\newif\ifqf@full
\newif\ifqf@boxed
\newsavebox{\qf@box}
\newcommand{\qf@open}{\global\qf@fulltrue\par\global\qf@boxedtrue
  \setbox\qf@box\vbox\bgroup\hsize\textwidth\linewidth\textwidth}
\newcommand{\qf@emit}[1]{\par\egroup\global\qf@boxedfalse
  \ifdim\dimexpr\ht\qf@box+\dp\qf@box\relax<#1\textheight\noindent\usebox{\qf@box}\par\else\unvbox\qf@box\fi}
\newcommand{\qfullsplit}{\ifqf@boxed\qf@emit{0.85}\fi}
\newcommand{\qf@close}{\ifqf@boxed\qf@emit{0.45}\fi\global\qf@fullfalse}
\newcommand{\qfullbegin}{\end{multicols}\qf@open}
\newcommand{\qfullend}{\qf@close\begin{multicols}{2}}
\newcommand{\qfullstart}{\qf@open}
\newcommand{\qfulldone}{\qf@close}
\newcommand{\qfig}[2][1]{\qf@check{#2}%
  \ifqf@full\qf@blk{#1}{#2}\else
  \ifnum\@nameuse{qf@wide@#2}=1 \expandafter\@firstoftwo\else\expandafter\@secondoftwo\fi
  {\end{multicols}\begin{center}\qf@img{#1}{\textwidth}{#2}\end{center}\begin{multicols}{2}}%
  {\qf@blk{#1}{#2}}\fi}
\newcommand{\qf@blk}[2]{\par\smallskip{\centering\qf@img{#1}{\linewidth}{#2}\par}\smallskip}
\newcommand{\qfigi}[2][1]{\qf@check{#2}\qf@img[valign=c]{#1}{\linewidth}{#2}}
\makeatother
\newcommand{\papersec}[1]{\par\bigskip{\centering\Large\bfseries PAPER-#1\par}\medskip}
\newcommand{\partsec}[2]{\par\medskip\noindent{\large\bfseries PART-#1: #2}\par\smallskip}
\newcommand{\sect}[1]{\par\smallskip\noindent\textbf{#1}\par}
\newcommand{\qn}[1]{\par\smallskip\noindent\textbf{#1.}~}
\newcommand{\paraq}[2]{\par\smallskip\noindent\textbf{Paragraph for Questions #1 and #2}\par}
\newenvironment{qoptions}{\begin{itemize}[leftmargin=1.8em,itemsep=1pt,topsep=2pt,label={}]}{\end{itemize}}
\newcommand{\opt}[2]{\item[(#1)] #2}
% figure registry (generated by score2_tools.py finalize)
\qfigdef{p1_phy_s1_q01_1}{Score-I - Test 01 (16 Jun 2021)-img/p1_phy_s1_q01_1.png}{111.1}{0}
\qfigdef{p1_phy_s1_q03_1}{Score-I - Test 01 (16 Jun 2021)-img/p1_phy_s1_q03_1.png}{84.2}{0}
\qfigdef{p1_phy_s1_q05_1}{Score-I - Test 01 (16 Jun 2021)-img/p1_phy_s1_q05_1.png}{120.5}{0}
\qfigdef{p1_phy_s2_q06_1}{Score-I - Test 01 (16 Jun 2021)-img/p1_phy_s2_q06_1.png}{70.7}{0}
\qfigdef{p1_phy_s2_q09_1}{Score-I - Test 01 (16 Jun 2021)-img/p1_phy_s2_q09_1.png}{83.4}{0}
\qfigdef{p1_phy_s2_q10_1}{Score-I - Test 01 (16 Jun 2021)-img/p1_phy_s2_q10_1.png}{175.6}{0}
\qfigdef{p1_phy_s2_q10_2}{Score-I - Test 01 (16 Jun 2021)-img/p1_phy_s2_q10_2.png}{92.7}{0}
\qfigdef{p1_phy_s2_q10_3}{Score-I - Test 01 (16 Jun 2021)-img/p1_phy_s2_q10_3.png}{92.7}{0}
\qfigdef{p1_phy_s2_q13_1}{Score-I - Test 01 (16 Jun 2021)-img/p1_phy_s2_q13_1.png}{71.1}{0}
\qfigdef{p1_phy_s3_q01_1}{Score-I - Test 01 (16 Jun 2021)-img/p1_phy_s3_q01_1.png}{62.4}{0}
\qfigdef{p1_phy_s3_q02_1}{Score-I - Test 01 (16 Jun 2021)-img/p1_phy_s3_q02_1.png}{92.7}{0}
\qfigdef{p1_phy_s3_q04_1}{Score-I - Test 01 (16 Jun 2021)-img/p1_phy_s3_q04_1.png}{110.2}{0}
\qfigdef{p1_chem_s1_q02_1}{Score-I - Test 01 (16 Jun 2021)-img/p1_chem_s1_q02_1.png}{140.9}{0}
\qfigdef{p1_chem_s1_q04_1}{Score-I - Test 01 (16 Jun 2021)-img/p1_chem_s1_q04_1.png}{67.8}{0}
\qfigdef{p1_chem_s1_q04_2}{Score-I - Test 01 (16 Jun 2021)-img/p1_chem_s1_q04_2.png}{137.9}{0}
\qfigdef{p1_chem_s1_q04_3}{Score-I - Test 01 (16 Jun 2021)-img/p1_chem_s1_q04_3.png}{137.9}{0}
\qfigdef{p1_chem_s1_q04_4}{Score-I - Test 01 (16 Jun 2021)-img/p1_chem_s1_q04_4.png}{137.9}{0}
\qfigdef{p1_chem_s2_q07_1}{Score-I - Test 01 (16 Jun 2021)-img/p1_chem_s2_q07_1.png}{59.4}{0}
\qfigdef{p1_chem_s2_q07_A}{Score-I - Test 01 (16 Jun 2021)-img/p1_chem_s2_q07_A.png}{71.8}{0}
\qfigdef{p1_chem_s2_q07_B}{Score-I - Test 01 (16 Jun 2021)-img/p1_chem_s2_q07_B.png}{31.5}{0}
\qfigdef{p1_chem_s2_q07_C}{Score-I - Test 01 (16 Jun 2021)-img/p1_chem_s2_q07_C.png}{41.6}{0}
\qfigdef{p1_chem_s2_q07_D}{Score-I - Test 01 (16 Jun 2021)-img/p1_chem_s2_q07_D.png}{44.7}{0}
\qfigdef{p1_chem_s2_q08_1}{Score-I - Test 01 (16 Jun 2021)-img/p1_chem_s2_q08_1.png}{127.5}{0}
\qfigdef{p1_chem_s2_q08_2}{Score-I - Test 01 (16 Jun 2021)-img/p1_chem_s2_q08_2.png}{125.9}{0}
\qfigdef{p1_chem_s2_q09_1}{Score-I - Test 01 (16 Jun 2021)-img/p1_chem_s2_q09_1.png}{104.7}{0}
\qfigdef{p1_chem_s2_q11_1}{Score-I - Test 01 (16 Jun 2021)-img/p1_chem_s2_q11_1.png}{160.5}{0}
\qfigdef{p1_chem_s3_q01_1}{Score-I - Test 01 (16 Jun 2021)-img/p1_chem_s3_q01_1.png}{79.6}{0}
\qfigdef{p1_chem_s3_q03_1}{Score-I - Test 01 (16 Jun 2021)-img/p1_chem_s3_q03_1.png}{374.4}{1}
\qfigdef{p2_phy_s1_q01_1}{Score-I - Test 01 (16 Jun 2021)-img/p2_phy_s1_q01_1.png}{74.8}{0}
\qfigdef{p2_phy_s1_q02_1}{Score-I - Test 01 (16 Jun 2021)-img/p2_phy_s1_q02_1.png}{139.2}{0}
\qfigdef{p2_phy_s1_q04_1}{Score-I - Test 01 (16 Jun 2021)-img/p2_phy_s1_q04_1.png}{120.2}{0}
\qfigdef{p2_phy_s2_q07_1}{Score-I - Test 01 (16 Jun 2021)-img/p2_phy_s2_q07_1.png}{52.8}{0}
\qfigdef{p2_phy_s2_q08_1}{Score-I - Test 01 (16 Jun 2021)-img/p2_phy_s2_q08_1.png}{77.9}{0}
\qfigdef{p2_phy_s2_q10_1}{Score-I - Test 01 (16 Jun 2021)-img/p2_phy_s2_q10_1.png}{136.5}{0}
\qfigdef{p2_phy_s2_q12_1}{Score-I - Test 01 (16 Jun 2021)-img/p2_phy_s2_q12_1.png}{126.1}{0}
\qfigdef{p2_phy_s2_q14_1}{Score-I - Test 01 (16 Jun 2021)-img/p2_phy_s2_q14_1.png}{289.7}{1}
\qfigdef{p2_phy_s2_q14_2}{Score-I - Test 01 (16 Jun 2021)-img/p2_phy_s2_q14_2.png}{100.4}{0}
\qfigdef{p2_phy_s3_para15_16_1}{Score-I - Test 01 (16 Jun 2021)-img/p2_phy_s3_para15_16_1.png}{133.5}{0}
\qfigdef{p2_phy_s3_para17_18_1}{Score-I - Test 01 (16 Jun 2021)-img/p2_phy_s3_para17_18_1.png}{209.5}{0}
\qfigdef{p2_chem_s1_q03_A}{Score-I - Test 01 (16 Jun 2021)-img/p2_chem_s1_q03_A.png}{173.0}{0}
\qfigdef{p2_chem_s1_q03_2}{Score-I - Test 01 (16 Jun 2021)-img/p2_chem_s1_q03_2.png}{187.0}{0}
\qfigdef{p2_chem_s1_q03_C}{Score-I - Test 01 (16 Jun 2021)-img/p2_chem_s1_q03_C.png}{189.4}{0}
\qfigdef{p2_chem_s1_q03_D}{Score-I - Test 01 (16 Jun 2021)-img/p2_chem_s1_q03_D.png}{174.1}{0}
\qfigdef{p2_chem_s1_q05_1}{Score-I - Test 01 (16 Jun 2021)-img/p2_chem_s1_q05_1.png}{65.9}{0}
\qfigdef{p2_chem_s2_q09_1}{Score-I - Test 01 (16 Jun 2021)-img/p2_chem_s2_q09_1.png}{233.9}{0}
\qfigdef{p2_chem_s2_q12_1}{Score-I - Test 01 (16 Jun 2021)-img/p2_chem_s2_q12_1.png}{109.1}{0}
\qfigdef{p2_chem_s2_q12_2}{Score-I - Test 01 (16 Jun 2021)-img/p2_chem_s2_q12_2.png}{139.0}{0}
\qfigdef{p2_chem_s3_q15_1}{Score-I - Test 01 (16 Jun 2021)-img/p2_chem_s3_q15_1.png}{129.2}{0}
\qfigdef{p2_chem_s3_q15_A}{Score-I - Test 01 (16 Jun 2021)-img/p2_chem_s3_q15_A.png}{67.2}{0}
\qfigdef{p2_chem_s3_q15_B}{Score-I - Test 01 (16 Jun 2021)-img/p2_chem_s3_q15_B.png}{72.2}{0}
\qfigdef{p2_chem_s3_q15_C}{Score-I - Test 01 (16 Jun 2021)-img/p2_chem_s3_q15_C.png}{83.5}{0}
\qfigdef{p2_chem_s3_q15_D}{Score-I - Test 01 (16 Jun 2021)-img/p2_chem_s3_q15_D.png}{70.3}{0}
\qfigdef{p2_chem_s3_q16_A}{Score-I - Test 01 (16 Jun 2021)-img/p2_chem_s3_q16_A.png}{109.7}{0}
\qfigdef{p2_chem_s3_q16_B}{Score-I - Test 01 (16 Jun 2021)-img/p2_chem_s3_q16_B.png}{121.5}{0}
\qfigdef{p2_chem_s3_q16_C}{Score-I - Test 01 (16 Jun 2021)-img/p2_chem_s3_q16_C.png}{124.8}{0}
\qfigdef{p2_chem_s3_q16_D}{Score-I - Test 01 (16 Jun 2021)-img/p2_chem_s3_q16_D.png}{110.6}{0}
\qfigdef{p2_math_s1_q02_1}{Score-I - Test 01 (16 Jun 2021)-img/p2_math_s1_q02_1.png}{147.3}{0}
\qfigdef{p2_math_s1_q02_2}{Score-I - Test 01 (16 Jun 2021)-img/p2_math_s1_q02_2.png}{31.6}{0}

\begin{document}
\begin{center}{\Large\bfseries JEE (Main+Advanced) Score-I --- Test 01 (16 Jun 2021)}\end{center}
\begin{multicols}{2}
\papersec{1}
\partsec{1}{PHYSICS}
\sect{Section-I(i)}
\qn{1} For an ideal transformer, ratio of $V_1$ \& $V_2$ is equal to
\qfig{p1_phy_s1_q01_1}
Where $L_1$ \& $L_2$ are self inductances of primary and secondary windings. (For an ideal transformer coefficient of mutual induction for pair of coils is $M=\sqrt{L_1L_2}$) :-
\begin{qoptions}
\opt{A}{$\dfrac{L_1}{L_2}$}
\opt{B}{$\sqrt{\dfrac{L_1}{L_2}}$}
\opt{C}{$\dfrac{L_2}{L_1}$}
\opt{D}{$\sqrt{\dfrac{L_2}{L_1}}$}
\end{qoptions}
\qn{2} According to Sommerfield model of atom an electron revolve around nucleus in an elliptical orbit while nucleus remain at rest. If minimum and maximum distance of electron from nucleus is $r_1$ \& $r_2$ then find distance of electron from nucleus when it is at a position where line joining electron \& nucleus is perpendicular to major axis :-
\begin{qoptions}
\opt{A}{$r_1+r_2$}
\opt{B}{$r_2-r_1$}
\opt{C}{$\dfrac{r_1r_2}{r_1+r_2}$}
\opt{D}{$\dfrac{2r_1r_2}{r_1+r_2}$}
\end{qoptions}
\qn{3} A small block of mass $m$ is rigidly attached at `P' to a ring of mass $m$ and radius $r$. The system is released from rest at $\theta=30^\circ$ and rolls without sliding. The speed of centre of mass of system when block reaches the bottom is
\qfig{p1_phy_s1_q03_1}
\begin{qoptions}
\opt{A}{$\sqrt{rg}$}
\opt{B}{$\dfrac{\sqrt{rg}}{2}$}
\opt{C}{$\sqrt{\dfrac{rg}{2}}$}
\opt{D}{none}
\end{qoptions}
\qn{4} One mole of an ideal monoatomic gas is enclosed in a chamber at 300 K. The gas undergoes a process in which pressure is proportional to the volume. At the end of the process, the volume of the gas is doubled. The change in the internal energy of the gas is [$R$ is gas constant]
\begin{qoptions}
\opt{A}{450 $R$}
\opt{B}{700 $R$}
\opt{C}{1350 $R$}
\opt{D}{data insufficient}
\end{qoptions}
\qn{5} For a conducting cylinder with variable thermal conductivity along the axis the magnitude of rate of heat flow if temperature difference across ends is $(T_1-T_2)$,
\qfig{p1_phy_s1_q05_1}
\begin{qoptions}
\opt{A}{$6\pi k_0R(T_1-T_2)$}
\opt{B}{$\dfrac{\pi k_0R}{3\ln2}(T_1-T_2)$}
\opt{C}{$\pi k_0R(T_1-T_2)$}
\opt{D}{$\dfrac{4\pi k_0R}{\ln2}(T_1-T_2)$}
\end{qoptions}
\sect{Section-I(ii)}
\qn{6} A small bob can slide downwards from point A to B along either of the two different curved surfaces shown in cross-section in the diagram. These possible trajectories are circular arcs in vertical plane, and they lie symmetrically about straight line AB. During the motion, bob does not leave contact with surface.
\qfig{p1_phy_s2_q06_1}
\begin{qoptions}
\opt{A}{If friction is neglected every where bob will take a smaller time through path I as compared to time through path II.}
\opt{B}{If friction is neglected every where bob will take a smaller time through path II as compared to time through path I.}
\opt{C}{If friction is neglected every where bob will reach point B with same speed through both the paths.}
\opt{D}{If friction is significant along both paths then bob will reach B with a smaller speed, if it follows path II, as compared to if it follows path I.}
\end{qoptions}
\qn{7} On a spherical planet, the refractive index of atmosphere, as a function of altitude $h$ above the surface varies as
\[n(h)=\dfrac{n_0}{1+\dfrac{h}{2R}}\]
Where $n_0$ is a constant. A laser beam directed horizontally but an arbitrary altitude follows a trajectory that circles the planet.
\begin{qoptions}
\opt{A}{Radius of planet has a magnitude $R\sqrt2$}
\opt{B}{Radius of planet has a magnitude $2R$}
\opt{C}{At a height $R\sqrt2$ from the surface of planet, refractive index is $\dfrac{n_0}{2}$}
\opt{D}{At a height $R$ from surface, refractive index is $\dfrac{2n_0}{3}$}
\end{qoptions}
\qn{8} Two equal masses $m$ are at a distance $\ell$ apart and interact via gravity. They are given proper tangential speed $v_0$, so that they both travel in a circular of radius $\dfrac{\ell}{2}$ around their CM. If one of the masses is grabbed and held at rest, then the closest distance the other mass comes to it is $r$. Then :
\begin{qoptions}
\opt{A}{$v_0=\sqrt{\dfrac{Gm}{2\ell}}$}
\opt{B}{$v_0=\sqrt{\dfrac{Gm}{\ell}}$}
\opt{C}{$r=\dfrac{\ell}{3}$}
\opt{D}{$r=3\ell$}
\end{qoptions}
\qn{9} A uniform stick with mass $m$ and length $\ell$ lies on a frictionless horizontal surface. It is pivoted at a point a distance $x$ from centre. A spring (at its relaxed length) with spring constant $k$ is attached to the far end of the stick, perpendicular to the stick, as shown. Stick is given a small kick, then:
\qfig{p1_phy_s2_q09_1}
\begin{qoptions}
\opt{A}{Value of $x$ which yields largest frequency of oscillation is $\dfrac{\ell}{3}$.}
\opt{B}{Value of $x$ which yields largest frequency of oscillation is $\dfrac{\ell}{6}$.}
\opt{C}{Angular frequency of oscillation (largest value) is $\dfrac{2k}{m}$.}
\opt{D}{Angular frequency of oscillation (largest value) is $\dfrac{4k}{m}$.}
\end{qoptions}
\qn{10} The refractive index of the medium within a certain region $x>0$, $y>0$, changes continuously with $y$. A thin light ray travelling in air in the $x$-direction strikes the medium at right angles and moves through the medium along a circular arc of radius $R$.
\qfig{p1_phy_s2_q10_1}
\begin{qoptions}
\opt{A}{Refractive index of medium varies with $y$ as \qfigi{p1_phy_s2_q10_2}}
\opt{B}{Refractive index of medium varies with $y$ as \qfigi{p1_phy_s2_q10_3}}
\opt{C}{If refractive index of medium can increase upto a value $n=2.5$, the maximum value of $y$ is $\dfrac{3R}{5}$}
\opt{D}{If refractive index of medium can increase upto a value $n=2.5$, the maximum value of $y$ is $5R$}
\end{qoptions}
\qn{11} Two waves given by equation $y_1=A\sin(\omega t-kx)$ and $y_2=4A\sin(2\omega t-4kx)$ are traveling in two different strings having same mass per unit length and same cross-sectional area. Tension in first string is four times tension in the second string.
\begin{qoptions}
\opt{A}{Ratio of average energy density $U_1/U_2$ is $1/64$}
\opt{B}{Ratio of average intensity $I_1/I_2$ is $1/32$}
\opt{C}{Ratio of average power $P_1/P_2$ is $1/16$}
\opt{D}{Ratio of velocity of wave $v_1/v_2$ is $1/2$}
\end{qoptions}
\qn{12} A metal block is placed in a room which is at $10^\circ$C for long time. Now it is heated by an electric heater of power 500 W till its temperature becomes $50^\circ$C. Its initial rate of rise of temperature is $2.5^\circ$C/sec. The heater is switched off and now a heater of 100W is required to maintain the temperature of the block at $50^\circ$C.(Assume Newtons Law of cooling to be valid)
\begin{qoptions}
\opt{A}{The heat capacity of the block is 200 J/$^\circ$C}
\opt{B}{The heat capacity of the block is 100 J/$^\circ$C}
\opt{C}{The rate of cooling of block at $50^\circ$C if the 100W heater is also switched off is $0.5^\circ$C/s}
\opt{D}{The heat radiated per second when the block was $30^\circ$C is 50 W}
\end{qoptions}
\qn{13} Figure shows a square current carrying coil of edge length $L$. The magnetic field on the coil is given by $\vec B=\dfrac{B_0y}{L}\hat i+\dfrac{B_0x}{L}\hat j$ where $B_0$ is a positive constant. ($A$ is area of coil)
\qfig{p1_phy_s2_q13_1}
\begin{qoptions}
\opt{A}{If coil is free to rotate about $x$ axis torque on the coil is given by $\dfrac12iAB_0\hat i$.}
\opt{B}{If coil is free to rotate about $y$ - axis torque on coil is given by $-\dfrac12iAB_0\hat j$.}
\opt{C}{Resultant force on coil is zero.}
\opt{D}{Equation for the torque $\vec\mu\times\vec B$ where $\mu$ is magnetic moment of coil is not valid on the coil if any of the side is fixed as axis.}
\end{qoptions}
\sect{Section-IV}
\qn{1} A ring of radius $r$ made of wire of density $\rho$ is rotated about a stationary vertical axis passing through its centre and perpendicular to the plane of the ring as shown in figure. Find the percentage change in tension for 3 percent change in angular velocity of rotation. Ignore gravity.
\qfig{p1_phy_s3_q01_1}
\qn{2} A U-tube having uniform cross-section but unequal arm length $l_1=100$ cm and $l_2=50$ cm has same liquid of density $\rho_1$ filled in it upto a height $h=30$ cm as shown in figure. Another liquid of density $\rho_2=2\rho_1$ is poured in arm A. Both liquids are immiscible. What length of the second liquid (in cm) should be poured in A so that second overtone of A is in unison with fundamental tone of B. (Neglect end correction)
\qfig{p1_phy_s3_q02_1}
\qn{3} A source of frequency 175 Hz is placed between a man and a wall at the same height of person's ear. Find the velocity $v$ (assume $v\ll c$) of the source with which it should approach the wall such that person will detect 3 beats per second. [Take velocity of sound $c=350$ m/s]
\qn{4} A sphere is floating in a liquid in just submerged position as shown. Find the ratio of magnitude of forces exerted by liquid on lower half of sphere to that exerted by liquid on upper half of sphere (Neglect atmospheric pressure)
\qfig{p1_phy_s3_q04_1}
\qn{5} One mole of an ideal monoatomic gas is taken from state A to state B through the process $P=\dfrac32T^{1/2}$. It is found that its temperature increases by 100 K in this process. Now it is taken from state B to C through a process for which internal energy is related to volume as $U=\dfrac12V^{1/2}$. The total work performed by the gas (in Joule) is $W$, if it is given that volume at B is 100 m$^3$ and at C it is 1600 m$^3$. [Use $R=8.3$ J/mol-K]. Fill $W/87$.
\partsec{2}{CHEMISTRY}
\sect{Section-I(i)}
\qn{1} Which one of the following dimeric form of compounds is planar ?
\begin{qoptions}
\opt{A}{$\mathrm{B_2H_6}$}
\opt{B}{$\mathrm{Al_2Cl_6}$}
\opt{C}{$\mathrm{Al_2Br_6}$}
\opt{D}{$\mathrm{I_2Cl_6}$}
\end{qoptions}
\qn{2} Borazine $+\ \mathrm{H_2O}\longrightarrow \mathrm{H_3BO_3}+\mathrm{A_{(g)}}+\mathrm{B_{(g)}}$

$\mathrm{B+N_2}\longrightarrow \mathrm{A}$;\quad $\mathrm{B+Cl_2}\longrightarrow \mathrm{C}$

$\mathrm{A_{(g)}+C_{(g)}}\longrightarrow \mathrm{D_{(s)}}$
\qfig{p1_chem_s1_q02_1}
\begin{center}
\begin{tabular}{cc}
Flask - I & Flask - II\\
$T=300$ K & $T=400$ K\\
$P=485$ mm Hg & $P=600$ mm Hg
\end{tabular}
\end{center}
Two flasks -- I \& II of equal volume containing Gas `A' and Gas `C', are connected by a narrow tube of negligible volume. The two gases were prevented from mixing by stopper fitted in connecting tube. For further detail of experiment refer to the given figure. What will be final pressure in each flask when valve connecting two tubes is opened. Assume ideal gas behaviour.
\begin{qoptions}
\opt{A}{40 mm Hg}
\opt{B}{60 mm Hg}
\opt{C}{20 mm Hg}
\opt{D}{10 mm Hg}
\end{qoptions}
\qn{3} A single walled carbon nano--tube can be regarded an a single sheet of graphite rolled into a cylinder. Number of C--atom in a single walled nano--tube having length 1.0 mm \& diameter 1.08 nm ?

[Given : Carbon - Carbon bond length in graphite = 141.5 pm]
\begin{qoptions}
\opt{A}{$0.65\times10^{8}$}
\opt{B}{$2.6\times10^{8}$}
\opt{C}{$3.9\times10^{8}$}
\opt{D}{$1.3\times10^{8}$}
\end{qoptions}
\qn{4} For the interaction $\mathrm{H+H}\longrightarrow \mathrm{H_2}$, \textbf{INCORRECT} is --
\begin{qoptions}
\opt{A}{$\Psi_{\mathrm{H}}$ vs r \qfigi{p1_chem_s1_q04_1}}
\opt{B}{$(\Psi_{\mathrm{bonding}})_{\mathrm{H_2}}$ vs r \qfigi{p1_chem_s1_q04_2}}
\opt{C}{$(\Psi_{\mathrm{antibonding}})_{\mathrm{H_2}}$ vs r \qfigi{p1_chem_s1_q04_3}}
\opt{D}{$(\Psi_{\mathrm{antibonding}})_{\mathrm{H_2}}$ vs r \qfigi{p1_chem_s1_q04_4}}
\end{qoptions}
\qn{5} Identify the \textbf{CORRECT} option.

(I) AgCl is not soluble in $\mathrm{H_2O}$ but soluble in KCN.

(II) Zn when reacts with very dil $\mathrm{HNO_3}$(6\%) gives $\mathrm{NH_4NO_3}$.

(III) $\mathrm{CuSO_4}$ reacts with KI to form only $\mathrm{Cu_2I_2}+\mathrm{K_2SO_4}$.

(IV) $\mathrm{CuSO_4{\cdot}5H_2O}\xrightarrow{1000^\circ\mathrm{C}}$ The gases liberated are $\mathrm{SO_2}$ and $\mathrm{O_2}$.
\begin{qoptions}
\opt{A}{All statements are correct}
\opt{B}{I, II and IV are correct}
\opt{C}{II and III are incorrect statements}
\opt{D}{Only I is incorrect.}
\end{qoptions}
\sect{Section-I(ii)}
\qn{6} Which of the following options is \textbf{CORRECT} ?
\begin{qoptions}
\opt{A}{Tailing of mercury related to test of Ozone.}
\opt{B}{Carbon suboxide ($\mathrm{C_3O_2}$) is linear molecule and prepared by dehydration of Malonic acid}
\opt{C}{$\mathrm{K_2O}$ and $\mathrm{K_2O_2}$ are diamagnetic but $\mathrm{KO_2}$ and $\mathrm{KO_3}$ are paramagnetic}
\opt{D}{$\mathrm{Mn_2O_7}$ is acidic oxide but MnO is amphoteric}
\end{qoptions}
\qn{7}
\qfig{p1_chem_s2_q07_1}

An organic compound on reaction with $\mathrm{O_3}$ followed by Zn and $\mathrm{H_2O}$ gives P. The organic compound can be ?
\begin{qoptions}
\opt{A}{\qfigi{p1_chem_s2_q07_A}}
\opt{B}{\qfigi{p1_chem_s2_q07_B}}
\opt{C}{\qfigi{p1_chem_s2_q07_C}}
\opt{D}{\qfigi{p1_chem_s2_q07_D}}
\end{qoptions}
\qn{8} Select correct statement for following reactions :

\textbf{Reaction : I}
\qfig{p1_chem_s2_q08_1}
\textbf{Reaction : II}
\qfig{p1_chem_s2_q08_2}
\begin{qoptions}
\opt{A}{Reaction I gives mixture of 2 enantiomers \& Reaction II gives mixture of two diastereomers.}
\opt{B}{Both reactions give mixture of 2 enantiomers.}
\opt{C}{Both reactions give mixture of 2 diastereomers.}
\opt{D}{Both reactions give 2 fraction of organic compounds.}
\end{qoptions}
\qn{9} $\mathrm{CH_3{-}CH{=}CH{-}CHO}$ (excess) $\xrightarrow[\Delta]{\mathrm{LDA}}$
\qfig{p1_chem_s2_q09_1}
\begin{qoptions}
\opt{A}{`A' is anti--aromatic.}
\opt{B}{`B' is aromatic.}
\opt{C}{`A' is having degree of unsaturation = 5.}
\opt{D}{In the compound `C', C--H bond length is more than C--C bond length.}
\end{qoptions}
\qn{10} Find \textbf{CORRECT} statements.
\begin{qoptions}
\opt{A}{Diamond is better conductor of heat than graphite.}
\opt{B}{Diamond is better conductor of electricity than Ag.}
\opt{C}{Diamond is electrical insulater.}
\opt{D}{Diamond gets converted to graphite at very high pressure \& temperature.}
\end{qoptions}
\qn{11} From the phase diagram of water and an aqueous solution containing non volatile solute, identify the \textbf{correct} options:
\qfig{p1_chem_s2_q11_1}
\begin{qoptions}
\opt{A}{At temperature $T_0$, vapour pressure of solid solvent and vapour pressure of liquid solvent will be same.}
\opt{B}{Order of vapour pressure ($P_0$, $P_1$, $P_2$) are $P_0>P_1>P_2$.}
\opt{C}{$P_0=P_2\,e^{\displaystyle\frac{\Delta H_{\mathrm{fusion}}}{R}\,\frac{[T_0-T_1]}{T_0T_1}}$}
\opt{D}{$P_1=P_2\,e^{\displaystyle\frac{\Delta H_{\mathrm{vap}}}{R}\,\frac{[T_0-T_1]}{T_0T_1}}$}
\end{qoptions}
\qn{12} Which of the following statements are correct ?
\begin{qoptions}
\opt{A}{$\dfrac{k_p}{k_a}=e^{\Delta E/RT}$ where $k_p$ = Rate constant for reaction with catalyst.

$K_a$ = Rate constant for reaction without catalyst

$\Delta E$ = Difference in activation energy for the two ways of reactions.}
\opt{B}{Alcohol acts as a negative catalyst for the conversion of $\mathrm{Na_2SO_3}\longrightarrow \mathrm{Na_2SO_4}$ .}
\opt{C}{TEL (tetraethyl lead) acts as a negative catalyst in the knocking of petrol.}
\opt{D}{Dilute $\mathrm{H_2SO_4}$ acts as a positive catalyst in the hydrolysis of esters to carboxylic acids and alcohols.}
\end{qoptions}
\qn{13} In which of the following cases, the synergic bonding takes place at the $\pi^*$-orbital of the ligand.
\begin{qoptions}
\opt{A}{$[\mathrm{PtCl_3(C_2H_4)}]^{-}$}
\opt{B}{$[\mathrm{Ni(PF_3)_4}]$}
\opt{C}{$[\mathrm{Cr(C_6H_6)_2}]$}
\opt{D}{$[\mathrm{Fe(\pi-C_5H_5)_2}]$}
\end{qoptions}
\sect{Section-IV}
\qn{1} The internal energy ($U$) of an ideal (10.7 mole) gas is plotted against volume for a cyclic process ABCDA, as shown in the figure.
\qfig{p1_chem_s3_q01_1}
The temperature of the gas at B and C are 500 K and 300 K, respectively. The heat absorbed by the gas (in kcal) in this cyclic process is :
\qn{2} $E^{\circ}_{Cell}$ for the cell

\begin{center}$\begin{array}{l}\mathrm{Pt(s)\,|\,Fe^{2+}(aq)\,|\,Fe^{3+}(aq)\,||}\\ \mathrm{MnO_4^{-}(aq),H^{+}(aq)\,|\,Mn^{2+}(aq)\,|\,Pt}\ \text{ is 0.788 Volt}\end{array}$\end{center}

\textbf{Given :}

$\Delta G^{\circ}_f\,\mathrm{H_2O}(\ell)=-230\,\mathrm{kJ\,mol^{-1}}$

$\Delta G^{\circ}_f\,\mathrm{Fe^{3+}(aq)}=-10.2\,\mathrm{kJ\,mol^{-1}}$

$\Delta G^{\circ}_f\,\mathrm{Mn^{2+}(aq)}=-229\,\mathrm{kJ\,mol^{-1}}$

$\Delta G^{\circ}_f\,\mathrm{Fe^{2+}(aq)}=-84\,\mathrm{kJ\,mol^{-1}}$

Magnitude of $\Delta G^{\circ}_f\,\mathrm{MnO_4^{-}(aq)}$ in kJ mol$^{-1}$ is $x\times10^{2}$ then `$x$' is-

\textit{(Give your answer as nearest integer)}
\qfullbegin
\qn{3} Three reactions are given below, write sum of total number of hydrogen atoms on all the carbon atoms which are connected directly by a single bond to benzylic carbon (carbon connected to benzene ring) in the products.
\qfig{p1_chem_s3_q03_1}
\qfullend
\qn{4} (a) From Meridional and facial isomer of $[Ma_3b_3]^{n\pm}$ on replacement of only one `a' by `b', the number of isomer of the product obtained are respectively $x$ and $y$.

(b) The number of corner shared per tetrahedron for 2D-silicate is $z$.

Then $x+y+z$ is-
\qn{5} An optically active substance `A' is decomposing into optically active substance `B' and `C' as
\par\noindent $\mathrm{A}\xrightarrow{K=10^{-3}\,\mathrm{min}^{-1}}2\mathrm{B}+\mathrm{C}$\par
The specific rotations of A, B and C are $+40^\circ$, $+10^\circ$ and $-30^\circ$ per mole, respectively. Initially `A' and `C' were present in $4:3$ mole ratio. Close to how many hours the sample becomes optically inactive ?

[Given : $\ell\mathrm{n}\,\dfrac{7}{5}=0.34,\ \ell\mathrm{n}\,\dfrac{20}{13}=0.42$ ]
\partsec{3}{MATHEMATICS}
\sect{Section-I(i)}
\qn{1} Let the function $f:\mathbb{R}\to\mathbb{R}$ be defined as $f(x)=\min.(x+2,\,4-2x,\,1+4x)$. The maximum value of $f(x)$ is equal to -
\begin{qoptions}
\opt{A}{$\dfrac13$}
\opt{B}{$\dfrac12$}
\opt{C}{$\dfrac23$}
\opt{D}{$\dfrac83$}
\end{qoptions}
\qn{2} Number of real roots of the equation $\sqrt{1+3x}+\sqrt{3+5x}+\sqrt{5+x}=\sqrt{3-x}+\sqrt{5-3x}+\sqrt{1-5x}$ is -
\begin{qoptions}
\opt{A}{1}
\opt{B}{2}
\opt{C}{3}
\opt{D}{infinite}
\end{qoptions}
\qn{3} In a sequence $\{a_n\}$, $a_1=2$, $a_{n+1}=1-\dfrac{1}{a_n}$ for $n\ge1$. Let $P_n$ be the product of its first $n$ terms, then the value of $P_{2017}$ is
\begin{qoptions}
\opt{A}{$\dfrac{-1}{2}$}
\opt{B}{2}
\opt{C}{$\dfrac12$}
\opt{D}{1}
\end{qoptions}
\qn{4} Given that `$a$' is a root of the equation $x^2-x-3=0$ then $\dfrac{a^3+1}{a^5-a^4-a^3+a^2}$ equals
\begin{qoptions}
\opt{A}{$\dfrac43$}
\opt{B}{$\dfrac34$}
\opt{C}{$\dfrac49$}
\opt{D}{$\dfrac94$}
\end{qoptions}
\qn{5} If $p,q,r\in\mathbb{R}$ satisfy $4p+3q+r=7$, then the least value of $(2p^2+q^2+r^2)$ is equal to -
\begin{qoptions}
\opt{A}{$\sqrt7$}
\opt{B}{$\dfrac{7}{\sqrt2}$}
\opt{C}{$\dfrac{\sqrt7}{2}$}
\opt{D}{$\dfrac{49}{18}$}
\end{qoptions}
\sect{Section-I(ii)}
\qn{6} Let us define system of vectors $\vec u_1=\vec v_1$ and $\vec u_k=\vec v_k-\mathrm{proj}_{\vec u_1}(\vec v_k)-\ldots-\mathrm{proj}_{\vec u_{k-1}}(\vec v_k),\ k\ge2$, where $\mathrm{proj}_{\vec a}(\vec b)=\dfrac{(\vec b\cdot\vec a)\vec a}{|\vec a|^2}$ and $\vec v_1=2\hat i+\hat j-\hat k,\ \vec v_2=\hat i+\hat j+\hat k,\ \vec v_3=3\hat i-2\hat j-\hat k$. Which of the following is correct ?
\begin{qoptions}
\opt{A}{$[\vec u_1\ \vec u_2\ \vec u_3]\ne0$}
\opt{B}{$\vec u_2\cdot\vec v_1\ne0$}
\opt{C}{$\vec u_3\cdot\vec v_1=0$}
\opt{D}{$\vec u_3\cdot\vec v_2\ne0$}
\end{qoptions}
\qn{7} For a $3\times3$ invertible matrix $A$ satisfying the characteristic equation $A^3+pA^2+qA-rI_3=0$, which of the following is/are true [Here $\mathrm{tr}(A)$ = trace of matrix $A$, $\det(A)$=determinant value of matrix $A$]
\begin{qoptions}
\opt{A}{$p=-\mathrm{tr}(A)$}
\opt{B}{$q=\dfrac{(\mathrm{tr}(A))^2-\mathrm{tr}(A^2)}{2}$}
\opt{C}{$r=\det(A)$}
\opt{D}{$A^{-1}=\dfrac{-2\,(A-\mathrm{tr}(A)I_3)}{(\mathrm{tr}(A))^2-\mathrm{tr}(A^2)}$}
\end{qoptions}
\qn{8} Let $L_1:\dfrac{x+3}{-4}=\dfrac{y-6}{3}=\dfrac z2$ and $L_2:\dfrac{x-2}{-4}=\dfrac{y+1}{1}=\dfrac{z-6}{1}$ be two lines in $\mathbb{R}^3$. Which of the following statement(s) is(are) correct ?
\begin{qoptions}
\opt{A}{$L_1,L_2$ are coplanar.}
\opt{B}{$L_1,L_2$ are skew lines.}
\opt{C}{shortest distance between $L_1$ and $L_2$ is 9.}
\opt{D}{$(\hat i-4\hat j+8\hat k)$ is a vector perpendicular to both $L_1$ and $L_2$.}
\end{qoptions}
\qn{9} If $(x+1)(x+2)\cdots\cdots(x+99)=A_0+A_1\cdot x+A_2\cdot x^2+\cdots+A_{99}\cdot x^{99}$ and $S=\dfrac{\displaystyle\sum_{k=1}^{99}(k.A_k)}{100!}=\displaystyle\sum_{i=1}^{99}\left(\dfrac1{r_i}\right),\left(r_i<r_j,\forall i<j\right)$ then value of $\displaystyle\sum_{k=1}^{99}\dfrac{1}{\sqrt{(r_k-1)}}$ is lying between
\begin{qoptions}
\opt{A}{15 and 19}
\opt{B}{15 and 17}
\opt{C}{14 and 17}
\opt{D}{17 and 19}
\end{qoptions}
\qn{10} Volume of parallelopiped formed by planes $2x+2y+z=4$, $2x+2y+z=1$, $2x-y+z=3$, $2x-y+z=6$, $3x+4y+12z=-1$, $3x+4y+12z=12$ is $\dfrac pq$ where $p$ and $q$ are positive coprime integers then :
\begin{qoptions}
\opt{A}{$p+q=20$}
\opt{B}{$p+q=57$}
\opt{C}{$p-q=6$}
\opt{D}{$p-q=5$}
\end{qoptions}
\qn{11} A line divides a plane into 2 regions. Two lines divide the plane into maximum 4 regions. If $L_n$ is the maximum number of regions divided by $n$ lines then which of the following is/are true :
\begin{qoptions}
\opt{A}{$L_{20}=211$}
\opt{B}{$L_{10}=56$}
\opt{C}{$L_{15}=121$}
\opt{D}{$L_{25}=326$}
\end{qoptions}
\qn{12} Let $A(z_1)$ be the point of intersection of curves $\arg(z-2+i)=\dfrac{3\pi}{4}$ and $\arg(z+i\sqrt3)=\dfrac\pi3$. $B(z_2)$ be the point on the curve $\arg(z+i\sqrt3)=\dfrac\pi3$ such that $|z_2-5|$ is minimum and $C(z_3)$ be the centre of circle $|z-5|=3$. [Note : $i^2=-1$].

Then which of the following statement(s) is (are) correct ?
\begin{qoptions}
\opt{A}{Distance between $A(z_1)$ and $B(z_2)$ is equal to 3.}
\opt{B}{Distance between $A(z_1)$ and $B(z_2)$ is equal to 4.}
\opt{C}{Area of triangle ABC is equal to $3\sqrt2$.}
\opt{D}{Area of triangle ABC is equal to $2\sqrt3$.}
\end{qoptions}
\qn{13} Given that the real numbers $s,t$ satisfy $19s^2+99s+1=0$, $t^2+99t+19=0$ and $st\ne1$. The absolute value of $\dfrac{st+4s+1}{t}$, is
\begin{qoptions}
\opt{A}{a prime number}
\opt{B}{divisible by 19}
\opt{C}{divisible by 5}
\opt{D}{divisible by 2}
\end{qoptions}
\sect{Section-IV}
\qn{1} Consider the system of equations $ax+y=b$, $bx+y=a$, $ax+by=ab$, where $a,b\in\{0,1,2,3,4\}$. The number of ordered pairs $(a,b)$ for which system is consistent is equal to
\qn{2} Let $A=\begin{bmatrix}4&a\\-1&b\end{bmatrix}$ and $B=\begin{bmatrix}1&-2\\1&4\end{bmatrix}$ be two matrices satisfying the relation,

$A^3+3A^2B+3AB^2+B^3=(A+B)^3$. If $(A+B)=\lambda(A^{-1}+B^{-1})$, then $\lambda$ is equal to
\qn{3} Let the focus of conic $\dfrac{(y+x)^2}{16}-\dfrac{(y-x)^2}{8}=1$ is at $(a,b)$ then $\dfrac{a^2+b^2}{4}$ is
\qn{4} Let $z$ and $\omega$ be complex numbers such that $z+\omega=i$ and $z^2+\omega^2=1$. If the area of triangle formed by $z$, $\omega$ and origin is $A$, then $2\sqrt{80A^2+1}$ is equal to
\qn{5} Let $f:\mathbb{R}\to\mathbb{R}$ be non-constant differentiable function and satisfies, $f(x)=x^2-\displaystyle\int_0^1(x+f(t))^2\,dt$. Then $f(4)$ is equal to
\papersec{2}
\partsec{1}{PHYSICS}
\sect{Section-I(i)}
\qn{1} Time constant of the given circuit is $\tau$. If the battery is replaced by an ac source having voltage $V=V_0\cos\omega t$, power factor of the circuit will be
\qfig{p2_phy_s1_q01_1}
\begin{qoptions}
\opt{A}{$\omega\tau$}
\opt{B}{$\dfrac{1}{\sqrt{1+(\omega\tau)^2}}$}
\opt{C}{$\sqrt{1-(\omega\tau)^2}$}
\opt{D}{None}
\end{qoptions}
\qn{2} Coherent light consisting of two wavelengths $\lambda_1=4500$ \AA\ \& $\lambda_2=6000$ \AA\ is sent through both slits of a Young's double slit apparatus (Source is placed symmetrically for both the slits)
\begin{qoptions}
\opt{A}{Central maxima for both wavelengths will coincide}
\opt{B}{The third order bright fringe of $\lambda_1$ will coincide with the fourth order bright fringe of $\lambda_2$}
\opt{C}{The third order dark fringe of $\lambda_2$ will coincide with fourth order dark fringe of $\lambda_1$}
\opt{D}{The fringes of wavelength $\lambda_1$ will be wider than the fringes of wavelength $\lambda_2$}
\end{qoptions}
\qn{3} An equiconvex lens of focal-length 10 cm (in air) and refractive index 3/2 is put at a small opening on a tube of length 1m fully filled with liquid of refractive index 4/3. A concave mirror of radius of curvature 20 cm is cut into two halves $m_1$ and $m_2$ and placed at the end of the tube. $m_1$ and $m_2$ are placed such that their principal axes AB and CD respectively are separated by 1mm each from the principle axes of the lens. A slit S placed in air illuminates the lens with light of frequency $7.5\times10^{14}$ Hz(in air). The light reflected from $m_1$ and $m_2$ forms interference pattern on the left end EF of the tube taking lens surface as screen. O is an opaque substance to cover the hole left by $m_1$ and $m_2$. The width of the fringes on EF.
\qfig{p2_phy_s1_q02_1}
\begin{qoptions}
\opt{A}{180$\mu$m}
\opt{B}{220$\mu$m}
\opt{C}{120$\mu$m}
\opt{D}{60$\mu$m}
\end{qoptions}
\qn{4} A closed loop of cross sectional area $10^{-2}\ \mathrm{m^2}$ which has inductance $L=10$ mH and negligible resistance is placed in a time varying magnetic field. The figure shows the variation of $B$ with time for the interval 4 sec. The field is perpendicular to the plane of the loop (given at $t=0$, $B=0$, $i=0$). The value of the maximum current induced in the loop is
\qfig{p2_phy_s1_q04_1}
\begin{qoptions}
\opt{A}{0.1 mA}
\opt{B}{10 mA}
\opt{C}{100 mA}
\opt{D}{data insufficient}
\end{qoptions}
\qn{5} A planet of small mass $m$ moves around sun of mass $M$ along an elliptical orbit such that its minimum and maximum distance from sun are $r$ and $R$ respectively. Its period of revolution will be
\begin{qoptions}
\opt{A}{$2\pi\sqrt{\dfrac{(r+R)^2}{6GM}}$}
\opt{B}{$2\pi\sqrt{\dfrac{(r+R)^3}{3GM}}$}
\opt{C}{$\pi\sqrt{\dfrac{(r+R)^3}{2GM}}$}
\opt{D}{$2\pi\sqrt{\dfrac{(r+R)^3}{GM}}$}
\end{qoptions}
\qn{6} A parallel plate capacitor is charged using a battery, and the battery is then removed. The plates of the capacitor are then brought closer together. Which of the following statements is false ?
\begin{qoptions}
\opt{A}{The electric field inside the capacitor remains the same}
\opt{B}{The capacitance of the capacitor increases}
\opt{C}{The potential difference between the plates of the capacitor decreases}
\opt{D}{The energy stored in the capacitor increases}
\end{qoptions}
\sect{Section-I(ii)}
\qn{7} The system of two weights with masses $m_1$ and $m_2$ are connected with weightless spring as shown. The system is resting on the support S. The support S is quickly removed. Find the accelerations $a_1$ \& $a_2$ of the weights $m_1$ \& $m_2$ respectively immediately after the support S is removed from below given options. Use: $m_1=m_2=m$
\qfig{p2_phy_s2_q07_1}
\begin{qoptions}
\opt{A}{$a_2=2g$}
\opt{B}{$a_1=0$}
\opt{C}{$a_2=g$}
\opt{D}{$a_1=g/2$}
\end{qoptions}
\qn{8} ABCD is a smooth fixed vertical circular track. A small particle is projected tangentially from the lower most point A of the track with speed $u=\sqrt{4gr}$ where $r$ is the radius of track. AC is the vertical diameter of track, then
\qfig{p2_phy_s2_q08_1}
\begin{qoptions}
\opt{A}{The particle leaves the track somewhere between B and C on the path of track}
\opt{B}{The particle cannot hit the point C when it is moving in air}
\opt{C}{The particle cannot cross the line AC when it moving in horizontal direction}
\opt{D}{The particle cannot cross the line BD while moving vertically}
\end{qoptions}
\qn{9} Two narrow cylindrical pipes A and B are filled with hydrogen and helium gases respectively at same temperature. Pipe A is open at both ends and pipe B is open at one end and closed at the other end. If frequency of the second overtone of pipe A is equal to the frequency of the first overtone of the pipe B,
\begin{qoptions}
\opt{A}{ratio of the velocities of sound waves in the two pipes $\dfrac{v_A}{v_B}=\sqrt2$}
\opt{B}{$\dfrac{v_A}{v_B}=\sqrt{\dfrac{42}{25}}$}
\opt{C}{ratio of their lengths $\dfrac{l_A}{l_B}=\sqrt{\dfrac{168}{25}}$}
\opt{D}{$\dfrac{l_A}{l_B}=2\sqrt2$}
\end{qoptions}
\qn{10} Figure shows a cylindrical vessel of height 84 cm filled upto the brim. Cylinder is open from top. There are four holes in the vessel as shows in figure. The liquid falling at maximum horizontal distance from the vessel would come from:
\qfig{p2_phy_s2_q10_1}
\begin{qoptions}
\opt{A}{hole 1}
\opt{B}{hole 2}
\opt{C}{hole 3}
\opt{D}{hole 4}
\end{qoptions}
\qn{11} Choose the \textbf{CORRECT} statements from the following.
\begin{qoptions}
\opt{A}{When an electron is emitted from a nucleus, then its neutron/proton ratio is decreased.}
\opt{B}{Number of protons in the nucleus is conserved during $\alpha$-decay.}
\opt{C}{In a hydrogen atom following Bohr model, the acceleration of electron in $n=2$ orbit is less than that in $n=1$ orbit.}
\opt{D}{Potential energy of an electron in an atom decreases with increasing principal quantum number.}
\end{qoptions}
\qn{12} The switch is closed at $t=0$ in the adjoining circuit. Which of the following is correct?
\qfig{p2_phy_s2_q12_1}
\begin{qoptions}
\opt{A}{The potential difference across BD at $t=0$ is $2E/3$}
\opt{B}{The potential difference across AB at $t=\infty$ is $E/2$}
\opt{C}{The potential difference across BD at $t=0$ is $E/2$}
\opt{D}{The potential difference across AB at $t=\infty$ is $2E/3$}
\end{qoptions}
\qn{13} Each of the following system begins moving upwards with a constant acceleration. Select these cases in which quantity will not change due to this upward acceleration :
\begin{qoptions}
\opt{A}{time period of simple pendulum.}
\opt{B}{fraction of floating body submerged in a liquid.}
\opt{C}{time period of a spring block system.}
\opt{D}{pressure on the base of a container containing liquid.}
\end{qoptions}
\qfullbegin
\qn{14} The diagram shows the variation of $V$ and $I$ in an AC circuit. The circuit only be a series RC or series RL or series LC or series RLC. Consider the four different combinations of $V$ and $I$ graphs. Pick the correct combination/combinations for each graph. Solid curves represent $I$ and broken curves represent $V$.
\qfig{p2_phy_s2_q14_1}
\qfullsplit
\qfig{p2_phy_s2_q14_2}
Take the angular frequency of the AC voltage source to be 100 rad/s (This is $\omega$)

(I) : $R=1\,\mathrm{k}\Omega$, $L=1$H and $C=100\,\mu$F

(II) : $R=1\,\Omega$ and $L=10^{-2}$H

(III) : $R=100\,\Omega$ and $C=10^{+2}\,\mu$F

(IV) : $L=1$H, $C=100\,\mu$F and $R=1\,\Omega$
\begin{qoptions}
\opt{A}{a $\to$ II ; d $\to$ I, IV}
\opt{B}{c $\to$ I ; d $\to$ none}
\opt{C}{b $\to$ III ; c $\to$ IV}
\opt{D}{a $\to$ II,III ; b $\to$ II,III}
\end{qoptions}
\qfullend
\sect{Section-I(iii)}
\paraq{15}{16}
Shells I, II, III are conducting spherical neutral shells of radius $R$, $2R$ and $3R$ respectively. Shell I and shell III are connected. A point charge $q_0$ is now kept at centre O. Now Shell II is also grounded. $q_d$ represents charge on outer surface of shell-II. $q_b$ represents charge on outer surface of shell-I.
\qfig{p2_phy_s3_para15_16_1}
\qn{15} Choose the correct option
\begin{qoptions}
\opt{A}{$q_d=-3q_b$ .}
\opt{B}{The total charge on shell II is zero.}
\opt{C}{$q_d=-q_b$}
\opt{D}{$q_d=+3q_b$}
\end{qoptions}
\qn{16} The charge on surface `b' (outer surface of Shell I) is
\begin{qoptions}
\opt{A}{$(2/11)\,q_0$}
\opt{B}{$(3/11)\,q_0$}
\opt{C}{$(1/11)\,q_0$}
\opt{D}{none of these}
\end{qoptions}
\paraq{17}{18}
Consider an object as the head of the shown arrow of height 2mm. Assume only paraxial rays to be incident. If the first reflection is to be taken on mirror then on interface to create an image at $(x,y)$.
\qfig{p2_phy_s3_para17_18_1}
\qn{17} Find $2x+y$ :
\begin{qoptions}
\opt{A}{2 mm}
\opt{B}{2 cm}
\opt{C}{4 mm}
\opt{D}{4 cm}
\end{qoptions}
\qn{18} The created final image is -
\begin{qoptions}
\opt{A}{Real, enlarged then object}
\opt{B}{Real, same size as object}
\opt{C}{Virtual, enlarged then object}
\opt{D}{Virtual, diminished wrt object}
\end{qoptions}
\partsec{2}{CHEMISTRY}
\sect{Section-I(i)}
\qn{1} The element gallium (Ga) has the second largest liquid range of any element, melting at 29.78$^\circ$C and boiling at 2204$^\circ$C at atmospheric pressure.

Assume that you construct a thermometer using, gallium as the fluid instead of mercury, and that you define the melting point of gallium as 0$^\circ$G and the boiling point of gallium as 1000$^\circ$G. What is the melting point of sodium chloride (801$^\circ$C) on the gallium scale ?
\begin{qoptions}
\opt{A}{801}
\opt{B}{368}
\opt{C}{771.2}
\opt{D}{354.7}
\end{qoptions}
\qn{2} Compound X ($\mathrm{C_7H_{13}Br}$) is a tertiary bromide. On treatment with potassium ethoxide in ethanol, X is converted into Y ($\mathrm{C_7H_{12}}$). Ozonolysis of Y gives Z as the only product. Z can also be obtained by the oxidation of heptan-2, 6-diol.

Which statement about Z is incorrect.
\begin{qoptions}
\opt{A}{Z can also be obtained by oxidation of Y with conc. hot $\mathrm{KMnO_4}$ solution.}
\opt{B}{Z gives negative 2, 4-DNP test.}
\opt{C}{Z gives positive iodoform test.}
\opt{D}{Z can be reduced to a diol with $\mathrm{NaBH_4}$}
\end{qoptions}
\qn{3} Which of the following disaccharides fulfills the two requirements stated below ?

(i) $\alpha$--(1,2) glycosidic linkage

(ii) Not a reducing sugar
\begin{qoptions}
\opt{A}{\qfigi{p2_chem_s1_q03_A}}
\opt{B}{\qfigi{p2_chem_s1_q03_2}}
\opt{C}{\qfigi{p2_chem_s1_q03_C}}
\opt{D}{\qfigi{p2_chem_s1_q03_D}}
\end{qoptions}
\qn{4} When an electron is removed from every ion of 1 mol of $\mathrm{H^-}$ ions, 72 kJ of energy is required. In another experiment, when electrons are removed from $\frac12$ mole of H (one electron each) and given them to another $\frac12$ mole of H (one electron each), 619.5 kJ of energy is found to be required. The ionisation energy (in kJ/mol) of hydrogen is :
\begin{qoptions}
\opt{A}{691.5}
\opt{B}{$-691.5$}
\opt{C}{1311}
\opt{D}{$-1311$}
\end{qoptions}
\qn{5} True statement(s) about this compound is/are:
\qfig{p2_chem_s1_q05_1}
\begin{qoptions}
\opt{A}{It is E isomer}
\opt{B}{It is Z isomer}
\opt{C}{It is optically active}
\opt{D}{It has center of symmetry}
\end{qoptions}
\qn{6} A molecule M associates in a given solvent according to the equation $\mathrm{M}\rightleftharpoons(\mathrm{M})_n$. For a certain concentration of M, the van't Hoff factor was found to be 0.9 and the fraction of associated molecules was 0.2. The value of $n$ is
\begin{qoptions}
\opt{A}{28/18}
\opt{B}{18/28}
\opt{C}{5}
\opt{D}{1/5}
\end{qoptions}
\sect{Section-I(ii)}
\qn{7} When excess of silver nitrate is reacted with $\mathrm{I_2}$ in aqueous solution then which of the following compounds are formed ?
\begin{qoptions}
\opt{A}{$\mathrm{AgIO_3}$}
\opt{B}{AgI}
\opt{C}{$\mathrm{HIO_3}$}
\opt{D}{$\mathrm{HNO_3}$}
\end{qoptions}
\qn{8} Which of the following is/are used as flux to remove impurities in metallurgical processes
\begin{qoptions}
\opt{A}{Silica}
\opt{B}{Limestone}
\opt{C}{NaCl}
\opt{D}{Coke}
\end{qoptions}
\qn{9}
\qfig{p2_chem_s2_q09_1}
\begin{qoptions}
\opt{A}{$\mathrm{II}\ \xrightarrow[\mathrm{(iii)\ Mg}]{\mathrm{(i)\ Ag_2O\ (ii)\ Br_2}}\ \xrightarrow{\ \mathrm{(iv)\ H_3O^+}\ }$ product is having boiling point more than cyclopropane.}
\opt{B}{VI gives optically active compound by reaction with $\mathrm{D_2}$(1 eq.) / Pt.}
\opt{C}{II is resolvable}
\opt{D}{Index of Hydrogen Defficiency of VI is 10}
\end{qoptions}
\qn{10} A Daniel Cell is made at 25$^\circ$C by connecting $\mathrm{Zn/Zn^{+2}}$ (0.1 M, 1 litre) \& $\mathrm{Cu/Cu^{+2}}$ (0.9 M, 1 litre) electrode. The cell is discharged till its emf reaches 1.10 volt, then it is charged (reversing the discharge process) by passing 0.6 F charge.

[\textbf{Given} : $E^{\circ}_{\mathrm{Zn/Zn^{+2}}}=0.76$ V, $E^{\circ}_{\mathrm{Cu/Cu^{+2}}}=-0.34$ V, $\dfrac{2.303RT}{F}=0.06$, $\log2=0.30$]

Select the correct option(s).
\begin{qoptions}
\opt{A}{Concentration of $\mathrm{Zn^{+2}}$ becomes 1 M when EMF reach 1.10 volt.}
\opt{B}{Concentration of $\mathrm{Cu^{+2}}$ becomes 0.5 M when EMF reach 1.10 volt.}
\opt{C}{EMF of the cell just after charging is 1.118 volt}
\opt{D}{After charging concentration of $\mathrm{Cu^{+2}}$ becomes 0.8 M.}
\end{qoptions}
\qn{11} Select incorrect statements :
\begin{qoptions}
\opt{A}{Pre--exponential factor for zero order reaction is a unitless quantity.}
\opt{B}{If $t_{1/4}=30$ sec then $t_{1/2}=60$ sec for first order reaction}
\opt{C}{If $t_{1/3}=30$ sec then $t_{2/3}=90$ sec for sec--order reaction}
\opt{D}{If $t_{1/5}=30$ sec then $t_{3/5}=90$ sec for zero order reaction.}
\end{qoptions}
\qn{12} \textbf{Correct} statements about indigotin is/are
\qfig{p2_chem_s2_q12_1}
\begin{qoptions}
\opt{A}{It shows geometrical isomerism}
\opt{B}{Trans is more stable than cis}
\opt{C}{Cis is more stable than trans}
\opt{D}{It is a planar compound}
\end{qoptions}
\qn{13}
\qfig{p2_chem_s2_q12_2}
Which of the following statements is/are correct about the above reaction ?
\begin{qoptions}
\opt{A}{Product shows geometrical isomerism.}
\opt{B}{Product shows colour due to extended conjugation.}
\opt{C}{Electrophile attack at para position due to large size.}
\opt{D}{Electron withdrawing group in phenol increases rate of reaction.}
\end{qoptions}
\qn{14} Which of the following are square planar complexes ?
\begin{qoptions}
\opt{A}{$[\mathrm{Ni(CO)_4}]$}
\opt{B}{$[\mathrm{PtCl_2(NH_3)_2}]$}
\opt{C}{$[\mathrm{Cu(NH_3)_4}]^{2+}$}
\opt{D}{$[\mathrm{FeCl_4}]^{-}$}
\end{qoptions}
\sect{Section-I(iii)}
\paraq{15}{16}
Ester having $\alpha$--`H' gives $\beta$-- keto ester in presence of alkoxide ion. This reaction is known as Claisen condensations reaction.
\qn{15}
\qfig{p2_chem_s3_q15_1}
Possible structure of \textbf{``Q''} is
\begin{qoptions}
\opt{A}{\qfigi{p2_chem_s3_q15_A}}
\opt{B}{\qfigi{p2_chem_s3_q15_B}}
\opt{C}{\qfigi{p2_chem_s3_q15_C}}
\opt{D}{\qfigi{p2_chem_s3_q15_D}}
\end{qoptions}
\qn{16} Which of the following compound is obtained during above interconversion.
\begin{qoptions}
\opt{A}{\qfigi{p2_chem_s3_q16_A}}
\opt{B}{\qfigi{p2_chem_s3_q16_B}}
\opt{C}{\qfigi{p2_chem_s3_q16_C}}
\opt{D}{\qfigi{p2_chem_s3_q16_D}}
\end{qoptions}
\paraq{17}{18}
An aqueous solution of salt \textbf{(P)} gives white precipitate with Barium chloride solution which is insoluble in dilute HCl \& dilute $\mathrm{HNO_3}$. When ammonia (aq) is added to the aqueous solution of salt \textbf{(P)} blue precipitate is obtained which dissolves in excess ammonia producing deep blue colouration due to formation complex \textbf{(Q)}. With KCN, \textbf{(P)} gives first a yellow precipitate which quickly decomposes to white precipitate along with liberating highly poisonous gas \textbf{(R)}. \textbf{(R)} on hydrogenation produces \textbf{(S)} which gives yellowish orange coloured complex \textbf{(T)} with $\mathrm{CoCl_3}$.
\qn{17} Which of the following options is \textbf{CORRECT} with respect to P, Q, R, S and T.
\begin{qoptions}
\opt{A}{`P' is $\mathrm{Ag_2SO_4}$, `Q' is $[\mathrm{Ag(NH_3)_2}]_2\mathrm{SO_4}$}
\opt{B}{`S' is a bidentate ligand \& `R' is cynogen}
\opt{C}{`R' is $\mathrm{CH_2N_2}$ \& `P' is $\mathrm{CuSO_4}$}
\opt{D}{None of the above options is correct}
\end{qoptions}
\qn{18} Which of the following options regarding complex Q \& T is \textbf{NOT CORRECT}.
\begin{qoptions}
\opt{A}{Q is paramagnetic where as T is diamagnetic.}
\opt{B}{Q has a square planar geometry where as T has octahedral geometry.}
\opt{C}{When dissolved in water, Van't Hoff coefficient for Q is 2 \& that of T is 4 assuming 100\% dissociation.}
\opt{D}{Only one of the above statements is correct.}
\end{qoptions}
\partsec{3}{MATHEMATICS}
\sect{Section-I(i)}
\qn{1} Let $a,b,c,d$ are roots of polynomial $x^4+5x^3-4x^2+6=0$ then the value of $\dfrac{1}{a^2-5a+6}+\dfrac{1}{b^2-5b+6}+\dfrac{1}{c^2-5c+6}+\dfrac{1}{d^2-5d+6}$ is :
\begin{qoptions}
\opt{A}{$\dfrac{-76}{46}$}
\opt{B}{$\dfrac{-219}{186}$}
\opt{C}{0}
\opt{D}{$\dfrac{677}{1426}$}
\end{qoptions}
\qn{2} A $2\times10$ rectangle needs to be tiled with tiles of size $1\times2$ (domino) as shown in figure. In how many ways can this be done?
\qfig{p2_math_s1_q02_1}
\qfig{p2_math_s1_q02_2}
\begin{qoptions}
\opt{A}{34}
\opt{B}{55}
\opt{C}{89}
\opt{D}{None of these}
\end{qoptions}
\qn{3} Number of pairs $(x,y)$ of real numbers that satisfy the equation
\[\left(\sin^2x+\frac{1}{\sin^2x}\right)^2+\left(\cos^2x+\frac{1}{\cos^2x}\right)^2=12+\frac12\sin y\]
\begin{qoptions}
\opt{A}{0}
\opt{B}{1}
\opt{C}{2}
\opt{D}{Infinite}
\end{qoptions}
\qn{4} Let $z_0,z_1,z_2,\ldots,z_6$ be distinct complex numbers satisfying $z^7-7=0$. If $p$ is the product of all sums of $z_i$ taken pairwise i.e. $p=(z_0+z_1)\cdot(z_0+z_2)\cdots\cdots(z_5+z_6)$, then $p^2$ is
\begin{qoptions}
\opt{A}{$7^3$}
\opt{B}{$7^6$}
\opt{C}{$2^{21}\cdot7^3$}
\opt{D}{None of these}
\end{qoptions}
\qn{5} $\dfrac{\displaystyle\sum_{n=1}^{99}\sqrt{10+\sqrt n}}{\displaystyle\sum_{n=1}^{99}\sqrt{10-\sqrt n}}$ equals
\begin{qoptions}
\opt{A}{$\sqrt2$}
\opt{B}{$\sqrt2-1$}
\opt{C}{$\sqrt2+1$}
\opt{D}{None of these}
\end{qoptions}
\qn{6} Let A, B, C, D be randomly selected subsets of $X=\{1,2,3,\ldots,n\}$ such that $A\cap C=\phi$, $B\cap D=\phi$. The probability of $A\cap B\ne\phi$, $B\cap C\ne\phi$, $C\cap D\ne\phi$ and $D\cap A\ne\phi$, is
\begin{qoptions}
\opt{A}{$(9^n-8^n+7^n-6^n+5^n)/\,9^n$}
\opt{B}{$(9^n-4.8^n+6.7^n-4.6^n+5^n)/\,9^n$}
\opt{C}{$(7^n-2.\,6^n+5^n)/\,7^n$}
\opt{D}{None of these}
\end{qoptions}
\sect{Section-I(ii)}
\qn{7} Let matrix $A=\begin{bmatrix}x&y&-z\\1&2&3\\1&1&2\end{bmatrix}$, where $x,y,z\in\mathbb{N}$ If $|\mathrm{adj}(\mathrm{adj}(\mathrm{adj}A))|=4^8\cdot5^{16}$, then the number of such matrix is $p$ and value of $\displaystyle\int_{-\pi/2}^{\pi/2}2^{\sin x}dx+\int_{5/2}^{4}\sin^{-1}(\log_2(x-2))dx$ is equal to $\dfrac{q\pi}{4}$. If $\vec a=\hat i+\hat j+2\hat k$, $\vec b=\hat i+2\hat j+2\hat k$ and $|\vec c|=1$ such that $[\vec a\times\vec b\ \ \vec b\times\vec c\ \ \vec c\times\vec a]$ has maximum value then
\begin{qoptions}
\opt{A}{the value of $|(\vec a\times\vec b)\times\vec c|^2$ is 0}
\opt{B}{$p=36$}
\opt{C}{$q=5$}
\opt{D}{$p-q=20$}
\end{qoptions}
\qn{8} For differential equation of \[y=c_1\left(e^{\sqrt x+c_2}+e^{-\sqrt x+c_2}\right)\] :
\begin{qoptions}
\opt{A}{Order = 1}
\opt{B}{Order = 2}
\opt{C}{Degree = 1}
\opt{D}{Degree = 2}
\end{qoptions}
\qn{9} Identify which of the following statement(s) is(are) correct.
\begin{qoptions}
\opt{A}{$\displaystyle\lim_{x\to0}\dfrac{x^{6000}-(\sin x)^{6000}}{x^2\cdot(\sin x)^{6000}}$ is equal to 1000}
\opt{B}{$\displaystyle\lim_{x\to0}\dfrac{x^{6000}-(\sin x)^{6000}}{x^2\cdot(\sin x)^{6000}}$ is equal to 3000}
\opt{C}{$\displaystyle\lim_{x\to0}\dfrac{1-\cos^5x\cos^32x\cos^33x}{x^2}$ is equal to 11}
\opt{D}{$\displaystyle\lim_{x\to0}\dfrac{1-\cos^5x\ \cos^32x\ \cos^33x}{x^2}$ is equal to 22}
\end{qoptions}
\qn{10} Find polynomials $f(x)$, $g(x)$ and $h(x)$ If they exist such that for all $x$,
\[|f(x)|-|g(x)|+h(x)=\begin{cases}-1&\text{if }x<-1\\3x+2&\text{if }-1\le x\le0\\-2x+2&\text{if }x>0\end{cases}\]
\begin{qoptions}
\opt{A}{$f(x)=3x+3$}
\opt{B}{$g(x)=\dfrac{5x}{2}$}
\opt{C}{$h(x)=-x+\dfrac12$}
\opt{D}{None of these}
\end{qoptions}
\qn{11} If equation $16x^4-mx^3+(2m+17)x^2-mx+16=0$ has four distinct roots forming a geometric progression, and one of the root is $\dfrac12$, then :
\begin{qoptions}
\opt{A}{common ratio of G.P. is 4}
\opt{B}{$m=170$}
\opt{C}{sum of the roots $=\dfrac{85}{8}$}
\opt{D}{Product of root = 3}
\end{qoptions}
\qn{12} A leaf is torn from a paperback novel. The sum of the numbers on the remaining printed pages is 15,000 :
\begin{qoptions}
\opt{A}{sum of page numbers of the torn leaf is 51}
\opt{B}{number of printed pages in the book is 174}
\opt{C}{number of printed pages in the book is 173}
\opt{D}{sum of page numbers of the torn leaf is 55}
\end{qoptions}
\qn{13} Pankaj and Prateek are Algebra genius. In a game Pankaj is going to pick three non-zero real numbers and Prateek is going to arrange the three numbers as the coefficients of a quadratic equation $ax^2+bx+c=0$. Pankaj wins the game if and only if the resulting equation has two distinct rational solutions else Prateek will win.
\begin{qoptions}
\opt{A}{Prateek has 100\% chance of winning}
\opt{B}{Pankaj has 100\% chance of winning}
\opt{C}{both have 50\% chance of winning}
\opt{D}{1 may be a root for winning strategy}
\end{qoptions}
\qn{14} For the equation $2(2^x-1)x^2+(2^{x^2}-2)x=2^{x+1}-2$ for real numbers $x$
\begin{qoptions}
\opt{A}{Number of real solutions = 4}
\opt{B}{Number of real solutions = 3}
\opt{C}{Sum of real solutions = 0}
\opt{D}{Sum of real solutions = 2}
\end{qoptions}
\sect{Section-I(iii)}
\paraq{15}{16}
The gamma function $y=\Gamma(x)$ is a continuous smooth function in practice that connects the points $(x,y)$ given by $y=(x-1)!$ at the positive integer values for $x$. It is uniquely defined by $\Gamma(x)=\displaystyle\int_0^\infty t^{x-1}e^{-t}dt$, for every real $x$ except non positive integers together with $\Gamma(1)=1$ and $\Gamma\!\left(\dfrac12\right)=\sqrt\pi$. Another important function is beta function $B(x,y)=\displaystyle\int_0^1t^{x-1}(1-t)^{y-1}dt=\dfrac{\Gamma(x)\Gamma(y)}{\Gamma(x+y)}$ for all $x,y\in\mathbb{R}^+$. Now answer the following questions :
\qn{15} $\displaystyle\int_0^{\pi/2}\sin^{10}x\cdot\cos^{12}x\,dx$
\begin{qoptions}
\opt{A}{$\dfrac{3^4\cdot5\cdot7^3}{2^{16}}\pi$}
\opt{B}{$\dfrac{3^4\cdot5\cdot7^4}{2^{16}}\pi$}
\opt{C}{$\dfrac{3^4\cdot5^2\cdot7^2}{2^{16}}\pi$}
\opt{D}{None of these}
\end{qoptions}
\qn{16} $\displaystyle\int_0^\infty t^be^{-at}dt=\quad(b\in\mathbb{N})$
\begin{qoptions}
\opt{A}{$\dfrac{\Gamma(b+1)}{a^{b+1}}$}
\opt{B}{$\dfrac{\Gamma(a)\Gamma(b)}{ab}$}
\opt{C}{$\dfrac{\Gamma(b+1)\Gamma(a+1)}{a^{b+1}}$}
\opt{D}{$\dfrac{\Gamma(a+1)}{a^{b+1}}$}
\end{qoptions}
\paraq{17}{18}
For all $x\ge1$, let $f(x)$ be the function defined as
\[f(x)=\begin{cases}[x]\cdot\left|x-[x]-\dfrac{1}{2[x]}\right|,&\text{if }x<[x]+\dfrac{1}{[x]}\\[2ex] f\!\left(x-\dfrac{1}{[x]}\right)&\text{otherwise}\end{cases}\]
(where $[\,.\,]$ denote greatest integer function)
\qn{17} The function $f(x)$ must satisfies the condition
\begin{qoptions}
\opt{A}{$f(4)=f\!\left(4\dfrac18\right)$}
\opt{B}{$f\!\left(4\dfrac14\right)=f\!\left(4\dfrac18\right)$}
\opt{C}{$f\!\left(4\dfrac13\right)=f(4)$}
\opt{D}{$f(4)=f\!\left(4\dfrac14\right)$}
\end{qoptions}
\qn{18} The number of points at which $f(x)$ is not differentiable in the interval $x\in(4,5)$ is
\begin{qoptions}
\opt{A}{8}
\opt{B}{7}
\opt{C}{4}
\opt{D}{2}
\end{qoptions}
\end{multicols}
\end{document}
